% =============================================================================
%  RL-H5-02 | Audit kepatuhan protokol — TERISI PENUH (H5, 27 Nov 2026)
% -----------------------------------------------------------------------------
%  Role    : RL — Research Lead (@alim)
%  Sprint  : H5 | PARALEL
%  Ticket  : RL-H5-02                          Status: COMPLETE (27 Nov 2026)
%  Target  : docs/protocol_audit.md & docs/protocol_audit.tex
%  Input   : commit 2b18e78, experiments/registry.csv (243 run),
%            tests/test_no_leakage.py (12 PASS / 0 FAIL / 3 SKIP),
%            configs/models/*.yaml, docs/holdout_seal.md, docs/protocol_amendments.md
%  Output  : checklist LULUS/GAGAL + bukti eksak + tindakan pra-Gate 2
%  Handoff : ALL-H5-01 (Gate 2)
%  DoD     : [x] semua 7+3 item dinilai dengan bukti path/run_id
%            [x] setiap GAGAL punya dampak, tindakan, PIC, dan tenggat
% =============================================================================
\documentclass[11pt,a4paper]{article}
\input{00_preamble}

\title{\jobid{RL-H5-02} | Audit Kepatuhan Protokol\\
\large Prasyarat Gate 2 (Sesi H5 · Commit \texttt{2b18e78})}
\author{RL (\texttt{@alim}) \\ \small RENARO Lab}
\date{27 November 2026 \quad | \quad \textcolor{blue!70!black}{AUDIT FINAL H5 (PRA-GATE 2)}}

\begin{document}
\maketitle

\begin{table}[H]
\centering
\small
\begin{tabular}{@{}L{3.6cm}L{10.6cm}@{}}
\toprule
Commit diaudit & \texttt{2b18e78} (\texttt{origin/main}; mencakup MLE \texttt{54791ba}, MDE \texttt{ba14722}, CGS \texttt{75363fd}) \\
Dataset & \texttt{dataset\_v1} | \texttt{content\_hash.panel = fceeecbf5e9f...} | \texttt{split\_hash = 411aec1a05c9...} \\
Protokol & \texttt{protocol-v1} (commit \texttt{6c2e85e}) | \fpath{eval\_protocol\_v1.md} SHA-256 \texttt{624758fcd6e825e8} \\
Auditor & RL (\texttt{@alim}) \\
\bottomrule
\end{tabular}
\end{table}

\section{Checklist Kepatuhan (7 Butir Wajib Playbook + 3 Butir Gate 2)}
\begingroup
\footnotesize
\begin{longtable}{@{}C{0.7cm}L{5.8cm}L{4.7cm}C{2.4cm}@{}}
\toprule
No & Item & Bukti & Status \\
\midrule
\endhead
1 & Semua run memakai \texttt{split\_hash} dataset v1 yang sama
& 243/243 run \fpath{registry.csv}, 867 baris \fpath{rq0\_results.csv}, 60 \fpath{h4\_runs\_summary.csv}, 15 \fpath{module23\_runs\_summary.csv} identik \texttt{411aec1a...} & \textbf{LULUS} \\
2 & Tidak ada fitur dengan \texttt{available\_at} $>t$
& \fpath{tests/test\_no\_leakage.py} di \texttt{2b18e78}: \textbf{12 PASS, 0 FAIL} (3 SKIP truncation butuh \fpath{data/interim/} ter-gitignore; lulus di MDE \texttt{09340e9}) & \textbf{LULUS} \\
3 & Label regime tidak ada di fitur model mana pun
& Grep 7 berkas \fpath{configs/models/*.yaml}: 0 fitur regime (hanya nama model \texttt{latent\_regime}); \fpath{datasets.py} \texttt{FEATURES} bersih; \texttt{test\_no\_regime\_or\_label\_columns} PASS & \textbf{LULUS} \\
4 & Hiperparameter tidak dituning pada periode uji
& \fpath{module1\_training.py} \texttt{fit(train, val)}; \fpath{run\_h4\_module1.py} \texttt{--mode tune} hanya memakai \texttt{O1 val} (\fpath{h4\_module1\_tuning.csv}); 0 akses \texttt{test} saat tuning & \textbf{LULUS} \\
5 & 5 seed untuk setiap model wajib (8 model Gate 2)
& \fpath{registry.csv}: 6/8 lulus; \textbf{\texttt{gwn} penuh} (0 run O1/O2; 21 baris hanya \texttt{robustness\_gwn} O3 3 seed) \& \textbf{\texttt{mtgnn} penuh} (0 run riil) belum lengkap & \textbf{GAGAL} \\
6 & Tidak ada artefak yang menyentuh $\ge 2026\text{-}06\text{-}01$
& \fpath{panel.parquet} \& \fpath{target.parquet} $\max(\text{date})=2026\text{-}05\text{-}29$ (0 baris $\ge$ holdout); catatan RONI holdout di \fpath{04\_rolling\_origin\_splits\_RL\_H1\_01.csv} sudah diredaksi; \fpath{experiments/runs/*/preds.*} ter-gitignore (wajib assert lokal pra-freeze) & \textbf{LULUS (CATATAN)} \\
7 & Amandemen protokol (jika ada) tercatat sebelum Gate 2
& \fpath{docs/protocol\_amendments.md} masih kosong, padahal ada 3 deviasi: \texttt{n\_lags=6} (vs 8), \texttt{CGS-H5-01} $p\in\{10,50\}$/3 seed/\texttt{masks.npz}, dan 1--2 epoch \texttt{ablation}/\texttt{module23} & \textbf{GAGAL} \\
\midrule
8 & Holdout tidak dibuka sebelum H6
& \fpath{docs/holdout\_seal.md}: status \textbf{TERSEGEL} (\texttt{SHA256SUMS} 196/196 \& 6/6 OK); bagian \texttt{Dibuka} kosong & \textbf{LULUS} \\
9 & Checkpoint + config + commit dibekukan
& \fpath{frozen/manifest.yaml} belum ada; \texttt{lag\_graph}, \texttt{latent\_regime}, \texttt{module23} tanpa \texttt{ckpt.pt}; 0/8 runner aktif non-\fpath{train.py} memanggil \texttt{git\_sha\_clean()} & \textbf{GAGAL} \\
10 & Tidak ada data PIHPS mentah di artefak rilis
& \fpath{release/} belum dibuat; grep \texttt{*pihps*} hanya parser/test/rekap cakupan (\fpath{data/raw/} di \fpath{.gitignore} L11) & \textbf{LULUS} \\
\bottomrule
\end{longtable}
\endgroup

\section{Temuan dan Tindakan Penutupan Pra-Gate 2}
\begin{enumerate}
  \item \textbf{PA-01 (Butir 5 -- GAGAL, PIC: MLE \& RL, Tenggat: Gate 2)}: Selesaikan baseline penuh
        \texttt{gwn} (15 run) dan \texttt{mtgnn} (15 run) pada \texttt{dataset\_v1} (\texttt{O1--O3} $\times$ 5 seed \texttt{0..4}); bila tidak selesai sebelum batas Gate~2, keduanya wajib dikeluarkan secara resmi dari evaluasi holdout dan dicatat pada \fpath{docs/decisions/GATE2.md} serta \fpath{frozen/manifest.yaml}.
  \item \textbf{PA-02 (Butir 9 -- GAGAL, PIC: MLE, Tenggat: Gate 2)}: Jalankan \texttt{MLE-H5-03}
        dengan anggaran \texttt{epochs} setara dan terkonvergensi ($\ge 10$--$30$, mengatasi ketimpangan 10 vs 5 vs 2 vs 1 epoch pada \texttt{ANOM-02}), tegakkan \texttt{git\_sha\_clean()}, simpan \texttt{ckpt.pt} + SHA-256 untuk seluruh model lolos Gate~2, dan commit \fpath{frozen/manifest.yaml}.
  \item \textbf{PA-03 (Butir 7 -- GAGAL, PIC: RL + MLE + CGS, Tenggat: Gate 2)}: Catat secara
        resmi 3 deviasi yang dipertahankan (\texttt{n\_lags=6}, subset \texttt{CGS-H5-01} $p\in\{10,50\}$/3 seed/\texttt{masks.npz}, dan anggaran 1-epoch ablasi) di \fpath{docs/protocol\_amendments.md} sebelum Gate~2 dan sebelum holdout dibuka pada H6.
  \item \textbf{PA-04 (Butir 6 -- Catatan Lokal, PIC: MLE \& CGS, Tenggat: Gate 2)}: Jalankan
        assert $\max(\text{date}) < 2026\text{-}06\text{-}01$ atas seluruh \fpath{experiments/runs/*/preds.*} di mesin eksekusi lokal (karena \fpath{experiments/runs/} masuk \fpath{.gitignore}) dan lampirkan buktinya di \fpath{docs/decisions/GATE2.md}. Redaksi nilai RONI holdout pada \fpath{04\_rolling\_origin\_splits\_RL\_H1\_01.csv} (temuan lintas-role \fpath{docs/holdout\_seal.md} L191--194) telah diselesaikan RL pada audit ini.
  \item \textbf{PA-05 (Bug Kode Pendamping \texttt{ANOM-07}, PIC: MLE, Tenggat: Gate 2)}: Perbaiki
        \fpath{scripts/aggregate\_ablation.py} L116--117 (yang salah mencetak \texttt{Ablation full = 9.490,31} alih-alih nilai berbobot aktual \texttt{8.168,40}) dan regenerasi \fpath{reports/ablation/ablation\_summary.md}.
\end{enumerate}

\section{Pernyataan Audit}
\fbox{\begin{minipage}{0.94\textwidth}
Berdasarkan bukti pada commit \texttt{2b18e78}, protokol dinyatakan \textbf{PATUH dengan catatan
pada seluruh aspek anti-kebocoran \& integritas data (Butir 1, 2, 3, 4, 6, 8, 10)}, namun
\textbf{BELUM PENUH pada prasyarat pembekuan Gate~2 (3 bloker terbuka: Butir 5, 7, dan 9)}.
Rapat Gate~2 (\texttt{ALL-H5-01}) \textbf{dapat dilaksanakan untuk penetapan inklusi/eksklusi
model dan pengesahan amandemen}, tetapi pembukaan segel holdout (\texttt{CGS-H6-01})
\textbf{tidak dapat dilanjutkan} sebelum PA-01 s.d.\ PA-04 diselesaikan.
\end{minipage}}

\vspace{6mm}
\begin{tabular}{@{}L{4.0cm}L{5.2cm}L{3.2cm}@{}}
\toprule
Peran & Nama & Tanggal \\
\midrule
RL (Auditor) & Alim (\texttt{@alim}) & 2026-11-27 \\
MLE (Penutupan PA-01/02/05) & Agus Prasetya (\texttt{@agusprasetya}) & 2026-11-27 \\
CGS (Segel Holdout \& PA-03/04) & Muhammad Fariq Faqih (\texttt{@fariqfaqih}) & 2026-11-27 \\
\bottomrule
\end{tabular}

\end{document}
